\capitulo{1}{Introducción}

\section{Contexto}
Las interacciones farmacológicas constituyen uno de los principales problemas asociados a la prescripción 
concomitante de múltiples medicamentos, ya que pueden modificar su comportamiento farmacocinético y 
farmacodinámico, alterando la eficacia terapéutica o incrementando el riesgo de aparición de efectos adversos \cite{rodrigues2016drug}.

Entre las interacciones farmacocinéticas, una de las más relevantes es la que se produce a nivel del metabolismo 
hepático, especialmente cuando varios fármacos son sustratos, inhibidores o inductores de una misma isoenzima del 
sistema del citocromo P450 (CYP450) \cite{flockhart2021table}. La administración simultánea de medicamentos metabolizados por una misma 
isoenzima puede dar lugar a fenómenos de competencia por el metabolismo, inhibición o inducción enzimática, 
modificando las concentraciones plasmáticas de uno o varios fármacos y aumentando el riesgo de toxicidad o de 
fracaso terapéutico \cite{denisov2015mechanism} \cite{guengerich2008cytochrome}. Las isoenzimas CYP3A4, CYP2D6, CYP2C19, CYP2C9 y CYP1A2 son las principales responsables del 
metabolismo de una gran proporción de los medicamentos de uso clínico y, por tanto, las más frecuentemente 
implicadas en este tipo de interacciones \cite{zanger2013cytochrome}.

La relevancia clínica de estas interacciones queda reflejada en numerosos casos descritos en la literatura. 
Entre los ejemplos más representativos se encuentran la administración concomitante de \textbf{simvastatina y claritromicina}, 
en la que la inhibición de la isoenzima CYP3A4 incrementa las concentraciones plasmáticas de la estatina y aumenta 
el riesgo de rabdomiólisis \cite{zhou2007clinically} \cite{deodhar2020mechanisms}; la combinación de \textbf{warfarina y amiodarona}, asociada a la inhibición de CYP2C9 y al 
consecuente aumento del efecto anticoagulante con mayor riesgo de hemorragia \cite{lee2024review} \cite{deodhar2020mechanisms}; y la administración conjunta de 
\textbf{omeprazol y clopidogrel}, donde la inhibición de CYP2C19 reduce la activación del clopidogrel y puede disminuir 
su efecto antiplaquetario \cite{lee2024review}.

En este contexto, la identificación temprana de las posibles interacciones mediadas por el sistema CYP450 resulta 
fundamental para optimizar la farmacoterapia, mejorar la seguridad del paciente y facilitar la toma de decisiones 
clínicas basadas en la evidencia.


\section{Motivación}
La mayoría de herramientas comerciales presentan varias limitaciones:
\begin{itemize}
\item No hacen explícito el proceso de decisión seguido para clasificar la gravedad de la interacción, limitándose a mostrar el resultado junto con una explicación del mecanismo farmacocinético.
\item No proporcionan alternativas terapéuticas justificadas incluyendo combinaciones consideradas seguras.
\item No integran de forma conjunta la información sobre posibles efectos adversos de las alternativas propuestas.
\end{itemize}
Por este motivo, se propone una herramienta capaz de detectar posibles interacciones entre los principios activos introducidos, hacer explícito el razonamiento seguido para estimar el nivel de riesgo a partir de las isoenzimas CYP implicadas, proponer alternativas terapéuticas consideradas de bajo riesgo por la aplicación e informar sobre los posibles efectos adversos asociados.

\section{Descripción del proyecto}
El proyecto consiste en el desarrollo de una herramienta informática orientada al análisis de posibles interacciones farmacológicas 
de los principios activos. Permite al usuario seleccionar dos principios activos y obtener una estimación 
del nivel de riesgo de interacción en función de las enzimas implicadas en su metabolismo, proporcionando además una explicación de los 
resultados obtenidos, información sobre los efectos adversos asociados y posibles alternativas terapéuticas basadas en la clasificación 
ATC.

Para ello se construye una base de conocimiento a partir de 3 bases de datos biomédicas, públicas de referencia: \textbf{DrugBank} \footnote{DrugBank: Base de datos en línea de acceso libre que contiene información detallada sobre fármacos, sus principios activos, mecanismos de acción, dianas terapéuticas e interacciones.}, Centro de Información Online de Medicamentos (\textbf{CIMA})\footnote{CIMA: Plataforma oficial de la Agencia Española de Medicamentos y Productos Sanitarios (AEMPS) que proporciona acceso a la información autorizada de los medicamentos comercializados en España, incluyendo fichas técnicas, prospectos y datos administrativos.} de la Agencia Española de Medicamentos y Productos Sanitarios (AEMPS) y \textit{Side Effect Resource} (\textbf{SIDER}) \footnote{\textit{Side Effect Resource} (SIDER): Base de datos pública que recopila información sobre medicamentos comercializados y las reacciones adversas descritas en su documentación oficial.}

A partir de estas bases de datos se realiza un proceso previo de extracción, transformación y normalización de la información, generando conjuntos de datos optimizados que permiten realizar consultas de forma eficiente durante la ejecución de la aplicación.

La solución desarrollada ha sido implementada en Python utilizando el framework Dash para la construcción de la interfaz web, mientras que la gestión y el procesamiento de los datos se realizan mediante la biblioteca Pandas. El resultado es una aplicación interactiva que facilita la consulta de información farmacológica de forma sencilla y comprensible, sirviendo como herramienta de apoyo para el estudio de posibles interacciones entre medicamentos.


\section{Estructura del documento}
La memoria del presente proyecto se encuentra estructurada en siete capítulos. Para empezar se introduce el contexto del trabajo, la motivación que justifica su realización, la descripción general del proyecto y la organización de la memoria. 

En el segundo capítulo se recogen los objetivos del proyecto, distinguendolos en tres; general, específicos y de aprendizaje. 

En el tercer capítulo se muestran los conceptos teóricos necesarios para comprender el trabajo realizado. En él se describen los procesos de metabolismo de los fármacos, las interacciones farmacológicas, el funcionamiento del sistema enzimático CYP450, el sistema de clasificación ATC y el estado del arte de las herramientas existentes para la detección de interacciones entre medicamentos. 

El cuarto capítulo describe la metodología seguida durante el desarrollo del proyecto. Se detallan las bases de datos biomédicas utilizadas, el proceso de extracción, integración y depuración de la información, las técnicas empleadas para resolver las ambigüedades y las herramientas de software utilizadas para la implementación de la aplicación.

En el quinto capítulo se presentan los resultados obtenidos, describiendo el funcionamiento de la herramienta desarrollada, las funcionalidades implementadas y los diferentes casos de prueba empleados.

El sexto capítulo expone la discusión de los resultados, analizando las principales aportaciones del proyecto, sus limitaciones y las implicaciones derivadas de la metodología empleada.

Finalmente, el séptimo capítulo recoge las conclusiones alcanzadas, valorando el grado de cumplimiento de los objetivos planteados y proponiendo posibles líneas de trabajo futuro para ampliar las capacidades de la herramienta desarrollada.
